% Zdroj: PDF priklady-jaro-2025.pdf, fyzická str. 26-27. Značka: specializace UI-SU. Zařazeno: S1 (bayesovské učení / pravděpodobnostní uvažování).
\section{Bayesovské učení}
\szzotazka{jaro 2025}{30}{Bayesovské učení}

\noindent\textbf{Zadání.}\par
Supermarket bere bedny jablek od tří dodavatelů, polovinu beden od A, třetinu od B a šestinu od C. Bedny od sebe nelze pohledem rozeznat. Jablka jednotlivých dodavatelů se mírně liší průměrem -- mají 9cm, 10cm, nebo 11cm. Vybrali jsme náhodně vzorky jablek od jednotlivých dodavatelů a spočítali jsme, jak běžné jsou různé velikosti jablek u různých dodavatelů. Výsledky tohoto měření, počty jablek různých velikostí rozdělené dle jednotlivých dodavatelů, jsou v tabulce níže.

\begin{center}
\begin{tabular}{c|ccc}
\toprule
Dodavatel & 9cm & 10cm & 11cm \\
\midrule
A & 5 & 4 & 1 \\
B & 9 & 8 & 3 \\
C & 1 & 8 & 1 \\
\bottomrule
\end{tabular}
\end{center}

\begin{enumerate}[nosep]
  \item Uveďte maximálně věrohodné odhady pravděpodobností velikosti jablka pro jednotlivé dodavatele.
  \item Definujte pojem \textit{maximálně věrohodná hypotéza}.
  \item Z téže bedny jsme náhodně vzali dvě jablka a změřili je -- jedno mělo průměr 10cm a druhé průměr 11cm. Která z hypotéz: ,,Jablka jsou od A.``, ,,Jablka jsou od B.``, ,,Jablka jsou od C`` je maximálně věrohodná? Odpověď zdůvodněte.
  \item Jaká je \textit{bayesovsky optimální predikce} pravděpodobnosti velikosti 10cm pro další jablko ze stejné bedny (při předchozích pozorováních 10cm a 11cm jablek z této bedny)? Svou odpověď zdůvodněte.
\end{enumerate}

\medskip
\noindent\textbf{Náčrt řešení.}\par
\begin{enumerate}[nosep]
  \item Velikosti máme v diskrétních intervalech. Maximálně věrohodné odhady diskrétního rozložení jsou podíl četností, tj. pravděpodobnosti jsou:
  \begin{center}
  \begin{tabular}{c|ccc}
  \toprule
  $P(\text{velikost} \mid h)$ & 9cm & 10cm & 11cm \\
  \midrule
  A & $.5$ & $.4$ & $.1$ \\
  B & $.45$ & $.4$ & $.15$ \\
  C & $.1$ & $.8$ & $.1$ \\
  \bottomrule
  \end{tabular}
  \end{center}
  \item (Definice pojmu maximálně věrohodná hypotéza -- hypotéza maximalizující věrohodnost dat $P(\text{data} \mid h)$.)
  \item Maximálně věrohodná hypotéza po pozorování 10 a 11 cm je $C$ dávající pravděpodobnost pozorování $8\%$ oproti $4\%$ $A$ a $6\%$ $B$, viz třetí sloupec v tabulce níže.
  \item Spočteme aposteriorní pravděpodobnost jednotlivých hypotéz a tou vážíme jejich predikce, predikujeme $50\%$ pro velikost 10 cm, výpočet viz tabulka:
  \begin{center}
  \begin{tabular}{c|c|c|c|c|c}
  \toprule
  $h$ & $P(h)$ & $P(\text{data} \mid h)$ & $P(h, \text{data})$ & $P(h \mid \text{data})$ & predikce $10\text{cm} * P(h \mid \text{data})$ \\
  \midrule
  A & $.5$   & $.04$ & $.02$   & $3/8$ & $.4 * 3/8$ \\
  B & $1/3$  & $.06$ & $.02$   & $3/8$ & $.4 * 3/8$ \\
  C & $1/6$  & $.08$ & $.04/3$ & $2/8$ & $.8 * 2/8$ \\
  \midrule
  & & \multicolumn{2}{c|}{$P(\text{data}) = .16/3$} & & $\Sigma = 0.5$ \\
  \bottomrule
  \end{tabular}
  \end{center}
\end{enumerate}
